\section{Limitations}\label{sec:limits}

\paragraph{Open mode exposes the key to a threshold.} A coalition of $t$ nodes takes the key, and more than $f$ Byzantine nodes can certify several digests under one nonce, which also yields the key. \Cref{thm:imposs} shows that the first is inherent. On account chains a residual balance stays under a key that the open committee has seen; an owner reduces this by sweeping to a fresh address.

\paragraph{Sealed mode bounds the loss and does not remove it.} Colluding nodes at the threshold can release the highest-fee variant, or release before the trigger. The loss is bounded by the fee ceiling and the intent. A released Bitcoin transaction stays valid after the window, so an owner who posts a second request must name the same pin or move the pin first.

\paragraph{Proof coverage.} All results are for static corruption. Resharing has no theorem, and its safety rests on the consistency of the board. Liveness is sketched. The realization of $\FOMS$ is open. Open ECDSA is proved in the generic group model with ratio resistance of the sighash as an assumption, and its bound degrades with the number of slots that a lagging honest node keeps live. \Cref{sec:status} lists the open items.

\paragraph{Verification coverage.} The state machine is checked by simulation at $n = 4$, $7$ and $9$ within small bounds, and bounded model checking returned no result. The Lean model is abstract and covers one slot of the vote rules.

\paragraph{Scope of chains and intents.} A chain must offer a conflict object for the pin. Addresses that are not the raw public key or its hash, and accounts whose keys rotate, are out of scope. The Solana adapter is an outline. The intent format carries payments only; a contract call needs a calldata field and a gas table entry. The builder specification names hash-to-curve encodings and replacement rules that must be checked against the standards of each chain before deployment.

\paragraph{Costs and operations.} Sealed message size grows with the product of variants, inputs and nodes, and resharing costs a quadratic number of bytes per polynomial. Intents, variants and nonce points are public, so the scheme offers no privacy for requests, and a posting fee is assumed against spam. The epoch keys of $t$ nodes, if compromised before erasure, open every dealing of that epoch. The owner's device must be honest and must not reuse a nonce across messages. The measurements are of a prototype in process: no consensus protocol, no disk synchronization of the slot ledger and no acceptance by a chain node is included.

\section{Conclusion}\label{sec:conclusion}

The owner of a key knows more than the key: the owner knows the line on which every signature lies, and can hand out the line one point or one slot at a time. Line dealing turns that into a one-message protocol for existing addresses. Sealed points give a spend limit that survives the corruption of every node, because the dealt objects are conflicting signatures and the chain admits one. Open lines give the committee the freedom to build the transaction, under a lock that precedes every share and with a limit that an honest threshold enforces. The impossibility result says that these are the two options, and the choice between them belongs to the owner. The committee runs no key generation ceremony and no presignature protocol, and the scheme uses no Paillier encryption, oblivious transfer or class group. Its cost is a few kilobytes from the owner and about \nineNodeMicrosLow{}\us{} of CPU per node per request at nine nodes. What remains is proof work: resharing, adaptive corruption and the realization of $\FOMS$.
